\chapter{The lifting problem: transport a closure and its target}
\label{tx:chap:contract}
The syntax of an NS or PC proof has now been replaced by a finite
membership calculation. We can therefore ask a more general linear
question before constructing any operator: which parts of that calculation
must survive a change of universe?

We separate three tasks: identify the actual polynomial maps, preserve
the subspace calculation, and detect whether a target remains outside
the consequence space. The finite constructions later in the book
will fulfil this single transport contract.

\section{From one eligible proof to an infinite progression}
It helps to see the whole argument before its linear-algebra interfaces.
Recall that $k$ bounds the labels needed by one variable, $s$ bounds
the labels in one axiom index, and $t$ bounds those in one target index.
We work over a fixed field of positive characteristic $p$.
Fix one proof degree $d\ge1$, and choose a support budget $r$ covering
the degree-$d$ data. A fully invariant target uses $r=kd+s$; the
general indexed-target calculation uses $r=2kd+s+t$. The completed
lifting theorem will say that, for
\[
 q=p^a>r,\qquad n\ge\max\{n_*,2r-1\},
\]
a prescribed target request has proofs of degree at most $d$ at size
$n$ exactly when it does at size $n+q$.
Here $n_*$ is a fixed starting rank after which the template and target
data are stable; its precise value is given in
Chapter~\ref{tx:ch:profiles}. The number $p$ is the characteristic,
and $q$ is the size of the added block, not necessarily the size of
the coefficient field.

The first implication can be repeated with the same $d,r,q$, because
every larger starting size still satisfies the displayed inequality.
Thus one proof gives a constant degree bound on an infinite
progression. The reverse implication, applied also at lower degree
thresholds, identifies exact recurring minimum degrees. This is the
bridge to the four profiles.

Now suppose a finite value $d$ cannot recur. It cannot occur at an
eligible starting size, so its occurrence forces
\[
                    n<\max\{n_*,2r-1\}.
\]
Since $r$ is affine in $d$, this is a linear lower bound on $d$ once
$n$ is beyond the fixed starting rank. In T4 a value does recur; the
two-way comparison and a bound on $q$ instead move an occurrence to
a representative of size $O(d)$, creating a finite linear peak.
The remaining mathematical task is to prove the displayed lifting
statement with its linear retained-label budget. The definitions below
split that task into explicit obligations.

\section{Label actions in linear-algebra language}\label{tx:sec:action-primer}
We need some language from linear algebra and representation theory.
It is useful language, but none of the words in this section should
stand for an unexplained theorem. We introduce them through the
polynomial spaces already familiar from Chapter~\ref{tx:ch:closure}.

\subsection{A vector space together with relabelling}
Consider the degree-one space
\[
                  M=Kx_1\oplus Kx_2\oplus Kx_3.
\]
The permutation $(1\ 2)$ exchanges $x_1$ and $x_2$ and fixes $x_3$.
It therefore sends
$a x_1+b x_2+c x_3$ to $b x_1+a x_2+c x_3$.
This is an invertible linear map on coefficient vectors. Other
permutations give other linear maps, and doing two relabellings in
succession has the same effect as their composite.

In general, a \emph{representation} of $S_n$ is a vector space $M$
equipped with these compatible linear actions of the label
permutations. A \emph{module} is another name for such a space here;
later the same action is extended to linear combinations of permutations.
The vectors remain ordinary coefficient vectors; the representation
records how relabelling acts on them.

For a finite labelled set $\mathcal U_n$, the notation $K[\mathcal U_n]$
means the vector space with one independent basis vector $[u]$ for
each $u\in\mathcal U_n$. Its vectors are
$\sum_{u\in\mathcal U_n}a_u[u]$. If permutations relabel the objects
$u$, they permute these basis vectors. This is a \emph{permutation
module}. Monomials of degree at most $d$ give precisely such a basis
for our polynomial ambient space.

An \emph{invariant subspace}, also called a \emph{submodule}, is a
subspace carried to itself by every label permutation. In the displayed
example, $K(x_1+x_2+x_3)$ is invariant, whereas $Kx_1$ is not.
The consequence spaces are invariant because relabelling an entire
legal proof gives another legal proof from the uniform axioms.
An individual proof line need not be invariant.

\subsection{Maps that respect the relabelling}
A linear map $f:M\to M'$ is \emph{equivariant} when
\[
                         f(\pi v)=\pi f(v)
\]
for every label permutation $\pi$. The equation says that it makes
no difference whether we relabel first or apply the map first.
For example, sending an axiom index $[a]$ to its polynomial $f_{a,n}$
is equivariant by uniformity. Ordinary variable multiplication is
equivariant when the \emph{same} permutation acts on the variable
and the polynomial being multiplied.

Equivariant maps encode rules that are compatible with renaming labels.
Their symmetry lets one common expression in permutations interact
with all the maps at once. Chapter~\ref{tx:chap:deletion} gives the
elementary calculation behind this compatibility.

\subsection{Kernel, image and quotient: three different questions}
For a linear map $f:A\to B$, its \emph{kernel} is
$\ker f=\{a:f(a)=0\}$; its \emph{image} is
$\im f=\{f(a):a\in A\}$. They record respectively what the map kills
and what the map can produce. If $f$ is equivariant, both are invariant
subspaces. Indeed, $f(a)=0$ implies $f(\pi a)=\pi f(a)=0$, and
$\pi f(a)=f(\pi a)$ is again an image vector.

For a subspace $H\subseteq W$, the quotient $W/H$ identifies two
vectors whose difference belongs to $H$. The quotient map is
\[
                  \rho_H:W\to W/H,\qquad g\mapsto[g]=g+H.
\]
Its kernel is exactly $H$. Thus
\[
                         [g]=0\text{ in }W/H
                              \quad\Longleftrightarrow\quad g\in H.
\]
This will express the question of whether a target is derivable.
The zero \emph{class} does not mean that the original polynomial is
zero: a nonzero derived polynomial has zero class modulo the
consequence space.

For example, let $W=Kx\oplus Ky$ and $H=K(x+y)$. Then
$[x+y]=0$, while $[x]\ne0$ and $[y]\ne0$, with $[x]=-[y]$.
To see that $[x]\ne0$, an identity $x=a(x+y)$ would force both
$a=1$ and $a=0$ by comparing the two coefficients.
This calculation works in characteristic two as well.
If $H$ is invariant, relabelling the representatives defines an action
on the quotient: $\pi[g]=[\pi g]$.
A \emph{subquotient} of $W$ is a quotient $U/H$ with
$H\subseteq U\subseteq W$ invariant. It records the directions in
$U$ that are still missing from $H$.

\subsection{What the word exact will mean}
Suppose $A\xrightarrow{f}B\xrightarrow{g}C$ are linear maps.
The sequence is \emph{exact at $B$} if $\im f=\ker g$:
the vectors arriving from $A$ are exactly those killed by the map to
$C$. This is a statement about two specified subspaces of $B$.
For instance,
\[
        0\longrightarrow H\longrightarrow W
                      \xrightarrow{\rho_H}W/H\longrightarrow0
\]
says that inclusion is injective, its image is the kernel of the
quotient map, and the quotient map is surjective.

Later we will construct an operation $\Del$ taking suitable spaces and
maps at rank $N=n+q$ to spaces and maps at rank $n$.
To say that this operation is \emph{exact} means that it preserves
these image--kernel relationships, with their actual inclusions.
It must carry an injection to an injection, and an onto map to an
onto map, on the specified class. It must also identify the transported
kernel and image \emph{inside the transported ambient space}.
Matching their dimensions would not do this.

The one-dimensional space $K$, with every permutation acting trivially,
is called the \emph{trivial module}. It is \emph{simple}: it has no
proper nonzero subspace. We will use this particular observation for
invariant targets, rather than assume a general structure theorem
about representations.

\subsection{Background for the transport construction}
There are four layers of background here. The first is ordinary
linear algebra: spans, linear maps, kernels and quotient spaces.
The second is representation theory: we keep track of relabelling
while doing that linear algebra. Since the coefficient field has
positive characteristic, this is \emph{modular} representation theory.
One familiar shortcut is unavailable: averaging over all permutations
would divide by $n!$, which is zero in the field once $n\ge p$.
The later construction supplies a different operator rather than
using this averaging argument.

A representation is \emph{semisimple} if it is a direct sum of simple
representations. Here \emph{simple} means that there is no proper
nonzero invariant subspace. In positive characteristic, an invariant subspace
need not have an invariant complementary subspace. For example, the
swap of two coordinates over a characteristic-two field fixes
$K(x_1+x_2)$. Any invariant complementary line would be an eigenline
for the swap, but its only eigenvalue is one and its fixed space is
precisely that same line. Thus no invariant complement exists.
This elementary example explains why a decomposition into invariant
pieces cannot simply be assumed in the later construction.

The third layer is the language of categories and functors: a way to
say that one operation changes many spaces and their connecting
maps consistently. We spell out that requirement in
Section~\ref{tx:sec:grammar}. The fourth is the language of exactness,
usually taught in homological algebra. Here we use just its
linear-algebra content: preserving kernels, images and quotient
relations. No knowledge of chain complexes or homology groups is
needed. For optional background on this terminology, see
\cite{StacksExact}; the required calculation is developed in this book.

The tensor product introduced next belongs to ordinary multilinear
algebra. In our application it has a particularly concrete meaning:
one register for each possible variable multiplication.

\section{The data at one universe size}
\label{tx:sec:multiplication-registers}
We now put names on the ingredients of one ordinary PC step.
\emph{Nothing in this section changes the universe size.}
In particular, the map $m_n$ below is ordinary polynomial
multiplication at one fixed rank; it is not the lifting operator.
We first understand this familiar operation, then ask how to transport it.
The letter $z$ in the worked polynomial example below is just a
polynomial variable. In Chapter~\ref{tx:chap:deletion}, $z$ will instead
denote a linear operator made from block permutations; that operator
will be defined there.

\subsection{Three spaces with three different jobs}
Fix a universe size $n$ and a proof-degree budget $d\ge1$.
The variable names are the members $v$ of the finite set
$\mathcal V_n$; write $x_v$ for the corresponding polynomial variable.
There may be many more variables than universe labels.
For example, variables $x_{ij}$ give $n^2$ names.
The two numbers must not be identified.

The spaces we shall use are
\[
 \begin{array}{ll}
 P_d(n)=P_{n,\le d}
   &\text{all polynomials allowed as proof lines},\\
 L_d(n)=P_{n,\le d-1}
   &\text{all polynomials small enough to multiply once},\\
 X(n)=K[\mathcal V_n]
   &\text{one independent basis vector $[v]$ per variable name}.
 \end{array}
\]
$P_d(n)$ contains both derivable and underivable polynomials.
$L_d(n)$ also contains both kinds; it imposes only a degree restriction.
If $V_j(n)$ is the space reached after $j$ closure stages, the part
we may actually multiply is $V_j(n)\cap L_d(n)$.
The space $X(n)$ records which variable we use. We do \emph{not}
require that variable itself to have a proof.

For an explicit coordinate picture, take variables $x,y,z$ and $d=2$.
Then
\[
 \begin{split}
 P_2(n)&=\Span\{1,x,y,z,x^2,xy,xz,y^2,yz,z^2\},\\
 L_2(n)&=\Span\{1,x,y,z\},\\
 X(n)&=\Span\{[x],[y],[z]\}.
 \end{split}
\]
Here the brackets indicate variable-name basis vectors, not quotient
classes. This is a fixed-instance illustration; it does not specify
a uniform axiom family or compute a lower bound.

\subsection{Multiplication by one specified variable}
For each variable $x_v$, consider the map
\[
     M_v:L_d(n)\to P_d(n),\qquad M_v(f)=x_vf.
\]
It is linear: $x_v(af+bg)=a\,x_vf+b\,x_vg$.
It stays inside $P_d(n)$ because $\deg f\le d-1$.
For $d=2$, multiplication by $z$ sends the coefficient vector
$a+bx+cy+ez$ to $az+bxz+cyz+ez^2$.
This is an ordinary matrix calculation.

Now we want one map that packages \emph{all} the maps $M_v$.
Its input should allow a separate polynomial $f_v$ for each variable:
one slot for what is multiplied by $x$, one for what is multiplied by
$y$, and so on. On these slots the output should be
$\sum_v x_vf_v$. A direct sum of one copy of $L_d(n)$ for each
variable does exactly this.

\subsection{Variable-labelled copies and tensor products}\label{tx:sec:tensor-registers}
The space of variable-labelled inputs is $X(n)\otimes L_d(n)$.
Give it one basis vector
$[v]\otimes\mu$ for each pair consisting of a variable name $v$
and a monomial $\mu$ of degree at most $d-1$.
Extend by linear combinations. Grouping these basis vectors by $v$
identifies the space with those variable-labelled copies of $L_d(n)$:
\[
       X(n)\otimes L_d(n)\ \simeq\
                     \bigoplus_{v\in\mathcal V_n}L_d(n).
\]
Every vector has a unique expression $\sum_v[v]\otimes f_v$.
The notation $[v]\otimes f$ means: put the polynomial $f$ in the
slot labelled by $v$, with zero in every other slot.

For the three-variable example this space has dimension $3\cdot4=12$.
It remembers inputs such as $[x]\otimes y$ and $[y]\otimes x$
as distinct vectors, even though their products are the same
polynomial $xy$. It also allows sums such as
$[x]\otimes f+[z]\otimes g$, with two different multiplication requests.

The tensor symbol is different from an ordinary Cartesian product.
Polynomial multiplication on pairs
$(a,f)\in X(n)\times L_d(n)$, after identifying $a$ with its
linear polynomial, is \emph{bilinear}: it is linear in either
argument when the other is fixed. It is not a linear map on the
vector space of pairs. For example,
\[
 ( [x],0)+(0,y)=([x],y),
\]
whereas multiplication is zero on the two summands and is $xy$
on their sum. The tensor product makes all these bilinear requests
into one \emph{linear} map. In a basis it is just a matrix on the
pair coordinates described above.

In particular,
$([x]+[y])\otimes f=[x]\otimes f+[y]\otimes f$ and
$[x]\otimes(f+g)=[x]\otimes f+[x]\otimes g$.
These rules distribute sums in either input; the monomial-pair
description already implements them.
The variable-labelled-copy interpretation uses the specified basis
of $X(n)$. It is not a claim that a general tensor product comes
with a preferred decomposition into copies.

The same coordinate construction applies to arbitrary finite-dimensional
spaces $A$ and $B$. After choosing bases $a_i$ and $b_j$, their tensor
product has basis $a_i\otimes b_j$. A bilinear map $\beta:A\times B\to C$
extends to a linear map on these coordinates by sending
$a_i\otimes b_j$ to $\beta(a_i,b_j)$. For linear maps
$f:A\to A'$ and $g:B\to B'$, the tensor map is defined by
$(f\otimes g)(a\otimes b)=f(a)\otimes g(b)$ and extended linearly.
Changing either basis changes the coordinates by this same rule;
the construction does not depend on which bases were chosen.
If labels act on both spaces, their \emph{diagonal action} is
$\pi(a\otimes b)=(\pi a)\otimes(\pi b)$. Each summand uses the same
permutation in both positions. This is the tensor action used when
multiplication and later detection pairings are transported.

\subsection{The single multiplication map}
Combining the variable multiplications gives the linear map
\[
 \boxed{\quad
 m_n:X(n)\otimes L_d(n)\longrightarrow P_d(n),\qquad
 m_n\!\left(\sum_v[v]\otimes f_v\right)=\sum_v x_vf_v.
 \quad}
\]
The subscript $n$ only tells us at which universe size the variables
and polynomials are being used. On a basis pair it sends
$[v]\otimes\mu$ to $x_v\mu$.
This determines the whole map by linearity.
Some later formulas write $x_v$ itself for the basis vector of $X(n)$;
at its output, $x_v$ always means the actual polynomial variable.

This map need not be injective. For instance,
$[x]\otimes y-[y]\otimes x$ is a nonzero input vector with output zero.
Different multiplication requests can produce the same polynomial;
their difference belongs to the kernel.
The map is not being asked to recover the request from its output.
We use its \emph{image} to collect the polynomials obtained by legal
multiplications.

\subsection{A complete closure step with cancellation}
Return to the three axioms
\[
 f_1=xy+x,\qquad f_2=xy+1,\qquad f_3=xz
                         \quad\text{over }\mathbb F_2,
\]
and keep the degree budget $d=2$. We first take all their linear
combinations, $V_0=\Span\{f_1,f_2,f_3\}$. Comparing the coefficients
of $xy$ and $xz$, as in Chapter~\ref{tx:ch:closure}, gives
\[
                         V_0\cap L_2=\Span\{x+1\}.
\]
Both $f_1$ and $f_2$ were quadratic. Their sum $x+1$ is linear
and is now eligible for multiplication.
The input $[z]\otimes(x+1)$ lies in
$X\otimes(V_0\cap L_2)$, and
\[
                 m_n([z]\otimes(x+1))=xz+z.
\]
Adjoining the old vector $f_3=xz$ gives the target $z$.
The corresponding PC lines are completely explicit:
\[
 \begin{array}{c|l|l}
 \text{line}&\text{polynomial}&\text{reason}\\ \hline
 1&xy+x&\text{input }f_1\\
 2&xy+1&\text{input }f_2\\
 3&xz&\text{input }f_3\\
 4&x+1&\text{add lines 1 and 2}\\
 5&xz+z&\text{multiply line 4 by }z\\
 6&z&\text{add lines 5 and 3}.
 \end{array}
\]
Every line has degree at most two. Multiplying line 1 or line 2
directly by $z$ would instead give a cubic and exceed the budget.
This is why the operation must first identify the low-degree
\emph{combinations}, rather than select low-degree inputs.

For a general current stage the same calculation reads
\[
 \underbrace{V_{j+1}(n)}_{\text{next stage}}
 =\underbrace{V_j(n)}_{\text{old consequences}}
  +\underbrace{m_n\bigl(X(n)\otimes
                   (V_j(n)\cap L_d(n))\bigr)}_{\text{allowed new multiples}}.
\]
The plus sign is a sum of subspaces: take all sums of one vector
from each displayed space. The image term is itself a subspace
because $m_n$ is linear.

\subsection{The input, inclusion and target maps}
There are two other PC maps besides multiplication.
Let $\mathcal A_{n,\le d}$ be the axiom indices whose polynomials
fit the degree budget. A vector in $K[\mathcal A_{n,\le d}]$
stores one coefficient for each such index. The map
\[
        \alpha_n:K[\mathcal A_{n,\le d}]\to P_d(n),
                  \qquad [a]\mapsto f_{a,n}
\]
forms the resulting linear combination of usable axioms.
Thus $V_0(n)=\im\alpha_n$. Distinct input columns may be dependent;
the image, not the number of columns, is the initial consequence space.
The map
$\ell_n:L_d(n)\hookrightarrow P_d(n)$ is simply inclusion:
the same polynomial is now regarded as a possible degree-$d$ line.
Remembering this inclusion is necessary when forming an intersection.

For NS, instead put one coordinate on each permitted pair $(\mu,a)$,
where $\mu$ is a monomial and
$\deg\mu+\deg f_{a,n}\le d$. Calling this set $\mathcal C_d(n)$,
the certificate map is
\[
 \phi_n:K[\mathcal C_d(n)]\to P_d(n),\qquad[\mu,a]\mapsto\mu f_{a,n}.
\]
Its image $I_d(n)$ is one span, with no PC iteration.
The set $\mathcal C_d(n)$ here consists of certificate \emph{indices};
it is not the representation class $\mathcal C_r(N)$ introduced
in Chapter~\ref{tx:chap:deletion}.

Finally, the requested targets have a separate column map
\[
       J_n:K[\mathcal B_n]\to P_d(n),\qquad[b]\mapsto G_{b,n}.
\]
We use this map into $P_d(n)$ only when every full target degree
is at most $d$. A budget below a stabilised target degree cannot
derive that target; we do not truncate the target to make it fit.
The map records what we ask to derive, not extra input axioms.
Adding its columns to $\alpha_n$ would change the problem.

\begin{checkpoint}
$P_d$ contains possible proof lines; $V_j$ contains the lines already
reachable at stage $j$; $L_d$ is the degree test; $X$ labels the permitted
variables. The map $m_n$ assembles ordinary variable multiples.
Only after understanding these fixed-rank ingredients do we ask for
an operation relating them at two different ranks.
\end{checkpoint}

\section{Building a consequence space by linear operations}\label{tx:sec:grammar}
We build each next consequence space by four familiar tasks:
find an intersection, make variable-labelled copies, apply the
multiplication matrix, and add the old space. We now spell out that
recipe and why kernels and images describe it.

\subsection{A direct sum stores two vectors; a sum combines them}
If $U,V$ are subspaces of the same ambient $P$, then $U\oplus V$
is the space of \emph{pairs} $(u,v)$. In contrast, $U+V\subseteq P$
is the space of outputs $u+v$. The linear map
\[
              \sigma:U\oplus V\to P,\qquad (u,v)\mapsto u+v
\]
has image $U+V$.
These spaces need not have the same dimension.
For $U=V=Kx$, the input has two independent coordinates,
while every output is a multiple of $x$. The pair $(x,-x)$
belongs to the kernel. The direct sum remembers the two requests;
the sum forgets how an output was divided between them.

\subsection{An intersection is found by solving a difference equation}
Keep the same ambient space $P$ and define
\[
              \delta:U\oplus V\to P,\qquad(u,v)\mapsto u-v.
\]
A pair belongs to its kernel exactly when $u=v$. The common vector
then belongs to both $U$ and $V$. Thus the map
\[
                  U\cap V\longrightarrow\ker\delta,
                              \qquad h\mapsto(h,h)
\]
is an isomorphism, with inverse $(u,u)\mapsto u$.
Strictly, the kernel is a space of pairs, whereas the intersection
is a subspace of $P$. The inverse just given, followed by inclusion
in $P$, specifies how to identify them.
This small distinction matters when the whole calculation is transported.

In the cancellation example, $U=V_0$ and $V=L_2$.
A kernel vector records a linear combination of the three axioms
which equals a degree-one polynomial. Solving the difference equation
gives the common line $K(x+1)$. It finds a cancelled combination,
not merely a subset of the original input columns.

\subsection{Read the closure recurrence as a four-step algorithm}
Suppose the current embedded space $V_j(n)\subseteq P_d(n)$ is known.
\begin{enumerate}
\item Find $U_j(n)=V_j(n)\cap L_d(n)$ by the difference-map kernel
      above. These are the already derived polynomials eligible
      for one further multiplication.
\item Make one copy of $U_j(n)$ for each variable name.
      This is $X(n)\otimes U_j(n)$, sitting in
      $X(n)\otimes L_d(n)$ by the actual inclusion $U_j(n)\subseteq L_d(n)$.
\item Apply the restriction of $m_n$ to that input space, and take
      its image. Call the image $M_j(n)$. It consists of all sums
      $\sum_v x_vf_v$ with $f_v\in U_j(n)$.
\item Apply the sum map from $V_j(n)\oplus M_j(n)$ to $P_d(n)$.
      Its image is the next space $V_{j+1}(n)$.
\end{enumerate}
All these operations are finite linear algebra at one rank.
Finding kernels or images is Gaussian elimination; making the
variable-labelled copies is just organising coordinates.
Each vector thereby obtained has a legal proof, by the constructive
argument in Proposition~\ref{tx:prop:closure}.

An unchanged dimension would not tell us that the right space had
been obtained. For example, $Kx$ and $Ky$ are isomorphic
one-dimensional subspaces of $Kx\oplus Ky$, but only the first
contains the target $x$. We must keep the embeddings and the
specified maps, not just abstract isomorphism types.

\subsection{Diagram and functor: the language behind the recipe}
A \emph{diagram} in this book is a collection of specified vector
spaces and linear maps between them. Calling the data of Section 5.3
the \emph{degree-$d$ polynomial diagram} simply says that we remember
the ambient, low-degree, variable, input and target spaces together
with their maps. The inclusions are part of the data.

The vocabulary comes from category theory. A \emph{category} organises
objects and the maps allowed between them, with identity maps and
composition. Here the objects are spaces with label actions and the
allowed maps are equivariant linear maps. A \emph{functor} coherently
assigns a new object to each object and a new map to each allowed
map: it preserves identities and composition. Thus, for maps
$A\xrightarrow{f}B\xrightarrow{g}C$,
\[
                \Del(gf)=\Del(g)\Del(f).
\]
This equation ensures that a composite calculation can be
transported either all at once or step by step.
The operation used here is linear on the allowed maps and must
also preserve the finite direct sums needed in the recipe.

An identification between transported and smaller-rank spaces must
also respect their connecting maps. For fixed-data maps $f_N$ and
$f_n$, write these identifications as $\eta_A$ and $\eta_B$. The
required square is
\[
\begin{CD}
\Del(A_N) @>{\Del(f_N)}>> \Del(B_N)\\
@V{\eta_A}VV @VV{\eta_B}V\\
A_n @>{f_n}>> B_n .
\end{CD}
\]
It \emph{commutes} when the two paths have the same value:
$\eta_B\Del(f_N)=f_n\eta_A$. A \emph{natural} identification satisfies
this compatibility for the allowed maps. This is the content of
naturality used here; it is not a claim that auxiliary witness
coefficients have a unique choice.

We have \emph{not} constructed this operation yet.
Write $\Del$ for the operation we want, from a larger rank
$N=n+q$ to the retained rank $n$. We must first identify the transported
fixed data with the data at $n$, so that the transported multiplication
map really is $m_n$, rather than some unrelated matrix of the same size.
We must then show that it preserves the kernels, images and sums
used by the four-step recipe.

If those statements hold, the induction is straightforward.
Transport the initial image $\im\alpha_N$ to $\im\alpha_n$.
If the current space has been transported, preserve the difference-map
kernel, the variable-labelled copies with their inclusion, the
multiplication image, and finally the sum image.
The next stage is therefore transported too.
At a stage past both stabilisation points, this gives transport of
the entire degree-$d$ consequence space.
No individual proof length has to be bounded or followed through the
change of universe.

This is why exactness is useful here. It is the preservation of a
concrete linear-algebra recipe. The formal contract below records the
three things which will make this reasoning valid; its existence
is the subject of the later construction.

\begin{definition}[Transport contract for a polynomial diagram]\label{tx:def:polynomial-diagram}
At a fixed degree $d\ge1$, a transport from rank $N$ to rank $n$
has three obligations.
\begin{enumerate}
\item \textbf{Identify the data.} Identify the transported ambient,
lower-degree and basis spaces with their smaller-rank counterparts,
carrying $\alpha,\ell,m,\phi,J$ to the actual smaller-rank maps.
For multiplication, the identification
$\Del(X(N)\otimes Y)\simeq X(n)\otimes\Del(Y)$ must respect the
inclusions of the subspaces $Y\subseteq L_d(N)$ used in the recurrence.
\item \textbf{Preserve the calculation.} Preserve kernels, images,
finite direct sums and their defining maps throughout the closure and
target calculation. This transports intersections, sums and every
closure stage as embedded subspaces.
\item \textbf{Reflect target vanishing.} For $H_N=I_d(N)$ or $C_d(N)$,
put $Y_N=(H_N+\im J_N)/H_N$. Require
$\Del(Y_N)=0\Longrightarrow Y_N=0$. This last condition says that
transport cannot hide a requested target outside the consequence space.
\end{enumerate}
\end{definition}

The next two lemmas separate the work: the first two obligations
transport the closure, while the third turns that transport into a
two-way target-membership test.
Throughout this discussion, transport includes the coherent maps and
finite direct sums just described. The tensor identification is
compatible with the actual subspace inclusions, not an independently
chosen isomorphism for each space.

\begin{lemma}[Transporting a closure recurrence]\label{tx:lem:closure-transport}
Suppose an operation $\Del$ takes the larger-rank ambient, lower-degree,
variable and usable-axiom spaces and maps to their actual smaller-rank
counterparts. Suppose it preserves kernels and images on all the spaces
used by the recurrence, and has the natural permutation-first-factor
identification
\[
 \Del(X(N)\otimes Y)\simeq X(n)\otimes\Del(Y).
\]
Then, as embedded subspaces of the identified smaller ambient space,
\[
                         \Del(V_j(N))=V_j(n)
                         \quad\text{for every }j.
\]
It also transports the stable PC closure. Transport of the certificate
map gives $\Del(I_d(N))=I_d(n)$.
\end{lemma}
\begin{proof}
The initial equality is preservation of the actual axiom-map image.
Given the equality at stage $j$, express the intersection as the kernel
of $(v,l)\mapsto v-l$. Preserve that kernel and its inclusion. Apply
the natural tensor identification and the actual multiplication map.
Preserve its image and then its sum with the old stage. This gives the
next equality. Both finite-dimensional chains eventually stabilise;
choose a stage past both stabilisation points to obtain the closure
statement. NS is the single certificate-map image calculation.
\end{proof}

\begin{checkpoint}
The content is compatibility of the diagrams, not equality of dimensions.
An isomorphism chosen separately for every space need not carry the
intersection, multiplication map or target vector to the right one.
\end{checkpoint}

\section{Preserving consequences and detecting missing targets}\label{tx:sec:target-reflection}
We now reach a separate question. Suppose the calculation of the
consequence space has been transported correctly. Does that already
tell us, in both directions, whether the prescribed target is derivable?
The answer depends on whether transport can erase the obstruction
to deriving it. We start with the ordinary quotient test.

\subsection{A target leaves a remainder modulo the consequences}
Let $H_N$ be $I_d(N)$ for NS or $C_d(N)$ for PC, inside
$W_N=P_d(N)$. For a single target $g$, its class $[g]\in W_N/H_N$
is zero exactly when $g$ has a proof within the degree budget.
Think of the quotient as recording only directions not yet proved.

For several targets the composite map
\[
 K[\mathcal B_N]\xrightarrow{J_N}W_N
                         \xrightarrow{\rho_{H_N}}W_N/H_N
\]
sends each target index to that target's unproved class.
Every target is derived exactly when this map is zero.
Its image is naturally the space
\[
                       Y_N=(H_N+\im J_N)/H_N.
\]
To check the formula, map $H_N+\im J_N$ into $W_N/H_N$.
Its kernel is $H_N$, and its image is the target classes just described.
Thus $Y_N$ is a space of missing target directions, not a new
polynomial or a new axiom.
As before, the full target degrees must fit $d$ before we use
$J_N$ into $P_d(N)$; there is no degree truncation.

For the small example $W=Kx\oplus Ky$, $H=K(x+y)$, a request for
$x+y$ has zero obstruction. A request for both $x$ and $y$ has
obstruction $W/H\ne0$. Knowing that their sum is a consequence
does not derive the two summands individually.

\subsection{What exact transport establishes, and what is still missing}
If all the actual maps are identified correctly, exactness takes the
target image and its quotient to the corresponding calculation at $n$:
\[
                             \Del(Y_N)\simeq Y_n.
\]
If $Y_N=0$, then $Y_n=0$. A derived request therefore remains derived
after this downward transport.
But from $Y_n=0$ we know only $\Del(Y_N)=0$.
To conclude $Y_N=0$ we need to know that this operation cannot
send a \emph{nonzero} missing-target space to zero.
That additional implication is called \emph{reflection of vanishing}.
In ordinary language: a missing consequence must still be detectable
after transport.

There is no conflict with exactness. Even if an injection is
transported to an injection, both its transported source and its
transported destination might be zero.
The map $0\to0$ is injective and surjective.
Preserving relations among surviving spaces does not by itself
require every original space to survive.

\subsection{An exact operation which really loses an obstruction}
Here is a concrete illustration. Work with pairs of vector spaces
$(A,B)$, with maps given by pairs of linear maps
$(f_A,f_B)$. Maps act separately on the two components and are not
allowed to mix them. Define $F(A,B)=A$: retain the first component
and forget the second.

Kernels, images, quotients and direct sums are computed componentwise.
For example, $\ker(f_A,f_B)=(\ker f_A,\ker f_B)$, and forgetting the
second component leaves $\ker f_A$. The same check works for images.
Consequently $F$ is an exact operation on these paired spaces.
This is not a claim that ordinary projection preserves every
intersection of arbitrary subspaces of $K^2$; our objects and allowed
maps are the componentwise ones just specified.

Now take
\[
 W=(K,K),\qquad H=(K,0),\qquad Q=(0,K),\qquad
                         J:Q\hookrightarrow W.
\]
The target map includes the second component. Its missing direction
has obstruction
$Y=(H+\im J)/H=(0,K)\ne0$.
After transport,
\[
             F(W)=K,\qquad F(H)=K,\qquad F(Q)=0,\qquad F(Y)=0.
\]
All the image--kernel calculations have been preserved, but the
target obstruction was forgotten.
This example explains the logical issue only; it is not the
polynomial deletion construction or a counterexample to our final
theorem. The additional survival lemma is precisely what prevents
this phenomenon on the relevant polynomial subquotients.

\subsection{How the actual proof will keep missing directions visible}
If a polynomial $g$ is outside $H$, elementary linear algebra supplies
a functional $\lambda$ with $\lambda(H)=0$ and $\lambda(g)=1$:
extend a basis of $H$ by $g$ and prescribe these values.
For $H=K(x+y)$, $g=x$, the example is
$\lambda(ax+by)=a-b$.
Such a functional certifies nonmembership without requiring an
invariant representative of a proof.

The permutation basis of the polynomial space has a dot product,
so a functional can be represented by a coefficient vector.
All functionals annihilating $H$ form $H^\perp$.
Concretely, this pairing is
$\langle a,b\rangle=\sum_\mu a_\mu b_\mu$ on monomial coordinates.
It is nondegenerate because each coordinate can be tested separately.
There is no positivity requirement over the coefficient field.
In characteristic two, for example, $(1,1)$ pairs to zero with itself
while pairing nontrivially with $(1,0)$.
Accordingly $H^\perp$ means the \emph{annihilator} of $H$ under this
pairing, not a chosen complementary subspace.
Permutations preserve the pairing because they permute its coordinates
simultaneously. The induced action on a functional is
$(\pi\lambda)(v)=\lambda(\pi^{-1}v)$; this is the usual dual action.
Pairing them with the missing classes gives a map onto $K$ whenever
the obstruction is nonzero. The pairing is compatible with
relabelings as a whole; an individually chosen separating functional
need not be invariant.

The formal argument in Chapter~\ref{tx:chap:assemble-transfer} uses
exactness and a tensor compatibility to show that, if deletion erased
the obstruction, it would erase the source of this onto pairing.
It would then produce a map from zero onto the surviving line
$\Del(K)\simeq K$, which is impossible.
This is the reason for a doubled ambient support budget.
One monomial coordinate uses at most $kd$ labels; a pair of such
coordinates uses at most $2kd$. We are bounding the coordinates
of the pairing's ambient basis, not the union of labels in an
arbitrary polynomial or functional.

We have described the idea here, not assumed the survival theorem.
The next lemma records the implication once survival is proved.

\begin{lemma}[The target-reflection contract]\label{tx:lem:target-contract}
In addition to the closure transport, suppose the actual target map is
transported and every nonzero subquotient of $P_d(N)$ survives deletion.
Then
\[
        \im J_N\subseteq H_N\quad\Longleftrightarrow\quad
        \im J_n\subseteq H_n.
\]
\end{lemma}
\begin{proof}
The two membership conditions are $Y_N=0$ and $Y_n=0$.
Exactness identifies $\Del(Y_N)$ with $Y_n$. Since $Y_N$ is a
subquotient of the ambient polynomial space, survival of nonzero
subquotients reflects its vanishing. This proves both directions.
\end{proof}

The survival lemma will come from a nondegenerate permutation-basis
pairing, with doubled support. It is proved in
Chapter~\ref{tx:chap:assemble-transfer}; it is not an extra unproved
premise of the final theorem.

\section{The easier invariant target}
The simplest important request is refutation: derive the polynomial
$1$. More generally, let $G_N$ be fixed by every label permutation.
The input space of its target map has only one coordinate, recording
the scalar coefficient with which that target is requested.
For a fully invariant polynomial $G_N$, the source is just $K$:
\[
                   j_N:K\to P_d(N),\qquad a\mapsto aG_N.
\]
Its image in $P_d(N)/H_N$ is zero or a copy of the simple trivial
module. The deletion of Chapter~\ref{tx:chap:deletion} fixes the trivial
module $K$, so exactness already reflects
whether this composite is zero. This explains why the invariant-target
budget is smaller than the general indexed-target budget. The constant
one is only one example of an invariant target.

Here is the argument without the module terminology. A map from a
one-dimensional space has either zero image or a one-dimensional
image. If the missing class $[G_N]$ is nonzero, its line has trivial
label action because $G_N$ was invariant. Deletion preserves the
image map and keeps that trivial line nonzero. It therefore cannot
hide this particular obstruction. An indexed request can instead
have a higher-dimensional missing space, so the doubled-support
pairing argument is needed there.

\section{The budget table}
The fixed variable support is at most $k$, axiom-index support at most
$s$, and target-index support at most $t$. The spaces have the following
safe bounds.
\begin{center}
\begin{tabular}{@{}ll@{}}
\toprule
Basis or ambient space & Support bound\\
\midrule
Degree-$d$ monomials & $kd$\\
Multiplier--axiom pairs & $kd+s$\\
Variable with degree-$d-1$ monomial & $k+k(d-1)=kd$\\
Target index & $t$\\
Two degree-$d$ monomial coordinates & $2kd$\\
\bottomrule
\end{tabular}
\end{center}
One affine choice covering the whole diagram and the survival argument is
\[
                           R(d)=2kd+s+t.
\]
For an invariant target the pairing is unnecessary, so use $R(d)=kd+s$.
Nothing in the PC recurrence increases the bound because its intermediate
spaces are subspaces of those fixed ambient spaces. In particular, the
support budget does not grow with proof length.

\section{The connection with representation stability}\label{tx:sec:stability-context}
The contract belongs to a familiar mathematical pattern: hold the
complexity of a labelled object fixed and let the number of available
labels grow. Here the complexity parameter is proof degree $d$, and
one monomial uses at most $kd$ labels. The ambient spaces grow, but
the allowed equality patterns among the labels in a basis object
come from a finite list. Representation stability studies how such
families of symmetric-group representations are related across ranks.
Three references help locate our argument within that setting.
The terminology in this section concerns the surrounding literature.
A \emph{categorical equivalence} matches objects and their allowed
maps, with an inverse up to compatible isomorphisms. An
\emph{abelian category} is a setting in which kernels, images and
quotients have the formal properties familiar from modules.
Spaces with equivariant linear maps are an example. Our construction
below will establish the needed maps directly.

\subsection{Church--Ellenberg--Farb: the broad fixed-degree picture}
An FI-module is a functor from finite sets and injections to vector
spaces: every injection has an assigned linear map, and compositions
are respected. Church--Ellenberg--Farb~\cite{CEF} turn representation
stability into a finite-generation condition. Their characteristic-zero
results give eventual character polynomials and polynomial dimensions;
fixed-degree configuration-space cohomology is one motivating example.
This provides a useful comparison with fixing polynomial degree here.
A character is the trace of a permutation acting on a representation.
On a permutation basis it counts fixed basis objects. For example,
the single-label space has character equal to the number of fixed
labels; the space on ordered pairs of distinct labels has character
$a(a-1)$ when a
permutation fixes $a$ labels. A \emph{character polynomial} describes
such traces using counts of cycles of specified lengths.
The cited cohomology examples are background applications of that
framework, not prerequisites for this book's proof.

Our consequence spaces need not form an FI-submodule under ordinary
inclusion. For example, the uniform axiom $f_n=\sum_{i<n}x_i$ has
degree-one consequence space $Kf_n$. Under $[n]\hookrightarrow[n+1]$,
ordinary inclusion takes it to $K\sum_{i<n}x_i$, which is not contained
in $Kf_{n+1}$. The new coefficient of $x_n$ proves the failure directly.
Thus fixed-degree FI intuition does not supply the consequence-transfer
theorem automatically. We have to construct the maps which preserve
the actual uniform equations.

\subsection{Harman: the close modular-periodicity antecedent}
Harman~\cite[Theorem 2.5 and Corollary 2.11]{Harman} constructs
periodic equivalences for bounded-weight permutation-module categories
and extends them to the relevant abelian categories. The stable range
is linear: his permutation result compares $n,m>2r$. In positive
characteristic, the suitable residue of the rank matters, rather than
one characteristic-zero eventual-stability statement.
The structural principle of linear-range modular periodicity therefore
predates this proof-complexity application.
His Proposition~2.7 includes compatibility with natural operations on
permutation modules; Proposition~2.12 includes induction, restriction,
tensor products and duality on the abelian envelope. Compatibility with
these operations is therefore also an antecedent. Here we construct
the actual polynomial diagrams independently and prove that their
closure operations and prescribed-target obstructions survive deletion,
with the support and rank bounds needed for the degree application.
The bounded parameter in this comparison controls the small
non-growing part of a permutation-module index. In the tuple
models used here, the corresponding concrete parameter is the
number of labelled positions $b\le r$.

Keep the period convention precise. This book uses a prime power
\emph{strictly greater} than the support budget. Harman's printed
$p^{\lceil\log_p r\rceil}$ needs care when $r$ is itself a power of
$p$; the strict choice avoids that boundary. For example,
$\binom52=0$ and $\binom72=1$ modulo two, although $5$ and $7$ agree
modulo two. Counts of degree $r$ can detect a digit beyond that
non-strict modulus. The present argument proves its strict-period
statement directly.

\subsection{Hone and coauthors: an explicit functor viewpoint}
Hone, Klein, Pauwels, Sherman, Yacobi and Zhang~\cite{HKPSYZ}
construct explicit periodic functors through semisimplification,
motivated by the Duflo--Serganova functor in Lie superalgebra theory.
Their symmetric-group result in Theorem 8.5.3 gives a closely related
categorical setting. It helps explain why an operation which removes
a prime-power block can have more structure than a numerical
coincidence of dimensions. Its hypotheses and restricted categories
matter; this is not global exactness on all representations.
In that literature, \emph{semisimplification} quotients a category by
maps invisible to trace tests: $f:A\to B$ is negligible when
$\operatorname{tr}(gf)=0$ for every allowed $g:B\to A$.
This is an operation on maps, not merely replacing each representation
by a list of its simple factors.
The Duflo--Serganova analogy starts with a parity-changing square-zero
operator on a graded representation and takes its kernel modulo image.
For the present proof the concrete quotient below is sufficient;
none of the superalgebra machinery is assumed.

The concrete operation here is
$\Del(M)=\ker z_M/(\ker z_M\cap\im z_M)$. Here $z_M$ is a linear
operator on a representation $M$, formed from permutations of the
added block; Chapter~\ref{tx:chap:deletion} gives its definition.
Its powers mean repeated composition of that operator. In characteristic two,
$z^2=0$ makes this look like the homology of a square-zero operator.
Homology generally fails to preserve injections, and the two-point
example in Chapter~\ref{tx:chap:deletion} shows that failure directly.
For a square-zero map $d:M\to M$, the relation $d^2=0$ implies
$\im d\subseteq\ker d$. In homological language, kernel vectors are
\emph{cycles}, image vectors are \emph{boundaries}, and
$\ker d/\im d$ is the \emph{homology}: the cycles remaining after
boundaries are identified with zero.
This vocabulary describes a quotient of two concrete subspaces.
Our characteristic-two deletion has this form. In odd characteristic,
$z^p=0$ need not give $\im z\subseteq\ker z$, so the intersection
in the displayed quotient remains essential.
Our common group-algebra projector proves exactness on the required
bounded-support class, with the actual coefficient maps. In odd
characteristic the quotient is defined exactly as displayed, not
silently replaced by $\ker z/\im z^{p-1}$.

The proof which follows is self-contained: it constructs the witness,
proves exactness and transports the whole NS or ordinary-PC diagram.
The references explain the mathematical setting and the antecedents;
they are not unproved inputs to the argument. The application requires
prescribed-target reflection, the separate degree-zero case and the
onset/period estimates used in Chapter~\ref{tx:ch:profiles}. No claim
that modular periodicity itself originates here is needed.

\section{How the rest of the proof fulfils the contract}
We add $q=p^a>R(d)$ labels and retain at least $2R(d)-1$.
The prime-power block defines a nilpotent operator $z$. On permutation
bases a quotient of its kernel deletes exactly the block labels, but
that quotient is not automatically exact on subspaces. One common
group-algebra construction produces an idempotent $e$ which models the quotient
on every needed space at once: $e^2=e$, so $e$ acts as a projection.
That common expression supplies exactness
and map compatibility. The only hard quantitative question is how many
retained labels are needed to build it.

In characteristic two, the preliminary construction uses $1+2+\cdots+r$ labels for a support
budget $r$. The improved construction reuses one pool and requires only
$\max_{b\le r}(2b-1)$. Once these bounds are proved, this chapter's
contract becomes the actual transfer theorem rather than a conditional
plan. The next chapter first shows exactly how that theorem yields the
four types, so we know why the improvement from a sum to a maximum matters.
